\documentclass[11pt, a4paper]{article}
\usepackage[a4paper, left=1.5cm, right=1.5cm, top=1.8cm, bottom=1.8cm]{geometry}
\usepackage{fontspec}
\usepackage[english, bidi=basic, provide=*]{babel}
\babelprovide[import, onchar=ids fonts]{english}
\babelprovide[import, onchar=ids fonts]{thai}
\babelfont{rm}{Noto Sans}
\babelfont[thai]{rm}{Noto Sans Thai}
\usepackage{amsmath, amssymb}
\usepackage{booktabs}
\usepackage{longtable}
\usepackage{array}
\usepackage[table]{xcolor}
\usepackage{enumitem}
\setlist[itemize]{label=-}
\usepackage{hyperref}

\definecolor{primary}{HTML}{1E3A8A}
\definecolor{accent}{HTML}{0F766E}
\definecolor{rowtint}{HTML}{F8FAFC}

\hypersetup{
    colorlinks=true,
    linkcolor=primary,
    urlcolor=accent
}

\title{\textbf{\color{primary}ทฤษฎีสารสนเทศโฮโลแกรมสัมพัทธภาพ (RHIT V9.0):\\สถาปัตยกรรมสนามโฮโลแกรม 5 มิติ แกนกลางโซลิตอนสสารมืด และเสียงสะท้อนของวัตถุความหนาแน่นสูงยิ่งยวด}}
\author{เอกสารต้นฉบับเพื่อการตรวจสอบมาตรฐาน (รุ่น v9.0.0 / โครงสร้าง SSOT v9.0.0)}
\date{23 พฤศจิกายน ค.ศ. 2026}

\begin{document}

\maketitle

\begin{abstract}
\noindent\textbf{ทฤษฎีสารสนเทศโฮโลแกรมสัมพัทธภาพ (Relativistic Holographic Information Theory: RHIT V9.0)} นำเสนอทฤษฎีสนามควอนตัมเชิงประสิทธิผลจากล่างขึ้นบน (Bottom-Up Effective Field Theory) บนโครงสร้างเรขาคณิต 5 มิติแบบโค้งงอ $S^2 \times \mathbb{R}_t \times [u_{\text{IR}}, 0]$ โดยต่อยอดจากฐานรากรุ่น V8.5 ที่ผ่านการตรวจสอบอย่างสมบูรณ์ (\emph{Zero-Contamination Core}) ทฤษฎี RHIT V9.0 คงความถูกต้องของค่าสังเกตการณ์พื้นฐานทั้ง 22 ค่า ครอบคลุมจักรวาลวิทยา แรงโน้มถ่วง สเปกตรัมมวลควาร์กและเลปตอน เมทริกซ์ผสม CKM/PMNS และการละเมิดสมมาตร CP ที่ความแม่นยำ 15 ตำแหน่งทศนิยมจากแหล่งความจริงหนึ่งเดียว (SSOT) ด้วยค่าไคสแควร์รวม $\chi^2 = 1.9421$ ($\text{dof} = 22 - 18 = 4 > 0$, ค่าเบี่ยงเบนสูงสุด $0.8567\sigma$)

\vspace{0.5em}
\noindent นอกจากนี้ รุ่น V9.0 ยังเติมเต็มเสาหลักที่ 3 โดยพิสูจน์ค่าคงที่อันตรกิริยาในตัวเองกำลังสี่ $\lambda_{\text{V9}} = 2.750176 \times 10^{-83}$ และรัศมีแกนกลางโซลิตอนในสุญญากาศ $R_{\text{TF,vac}} = 320.3682\text{ pc}$ จากตัวแปรเรขาคณิต 5 มิติ ($u_{\text{IR}}, \delta_H$) โดยตรง เมื่อผสานกับกลไกปรับขนาดแกนกลางตามความหนาแน่น $T^\mu_{\ \mu}$ แบบ Symmetron แบบจำลองสามารถลดค่าไคสแควร์จลนศาสตร์เทียบกับฐานข้อมูล Local Volume Database (LVDB) สำหรับกาแล็กซีบริวารทางช้างเผือก ($N = 46$) ลงได้ $40.66\%$ และกลุ่มกาแล็กซีแคระท้องถิ่น ($N = 94$) ลงได้ $6.30\%$ เมื่อเทียบกับโปรไฟล์ NFW ของ $\Lambda\text{CDM}$ พร้อมวางระบบคำทำนายทางดาราศาสตร์ฟิสิกส์แบบทางเดียวสำหรับการสั่นพ้องของเมฆโบซอนรอบหลุมดำ M87$^*$ ($\alpha_g = 1.377$) และคลื่นความโน้มถ่วงสะท้อนหลังการรวมตัว
\end{abstract}

\section{สถาปัตยกรรมทางทฤษฎีและสมการแอ็กชันหลัก}

RHIT V9.0 ผสานเอนโทรปีข้อมูลควอนตัมบนขอบฟ้าทรงกลม $S^2$ การรวมแรง $SO(10) \to SU(4)_C \times SU(2)_L \times SU(2)_R \to \text{SM}$ ใน 5 มิติ แรงโน้มถ่วงอุบัติใหม่ และพลวัตของสสารมืดสเกลาร์เข้าไว้ในสมการแอ็กชันหลักหนึ่งเดียว:

\begin{equation}
S_{\text{RHIT}}^{\text{V9.0}} = S_{\text{Backend}}^{\text{5D RGE+GW}} + S_{\text{Gravity}}^{\text{Emergent+Starobinsky}} + S_{\text{Dark}}^{\text{Symmetron+V9}} + S_{\text{Matter}}^{\text{5D Complex SVD}}
\end{equation}

\begin{equation}
S_{\text{RHIT}}^{\text{V9.0}} = \int d^4x \sqrt{-g} \left[ \frac{c^3}{16\pi G_N(u_{\text{IR}})} \left( R + \frac{R^2}{6 M_{\text{sc}}^2} - \frac{6\gamma_{\text{ent}} H^2}{c^2} \right) - \frac{1}{2}(\partial\phi)^2 - V_{\text{sym}}(\phi) - \frac{\phi^2}{2M_s^2} T^\mu_{\ \mu} + \mathcal{L}_{5\text{D Matter}} \right]
\end{equation}

\subsection{ธรรมาภิบาลเชิงระเบียบวิธีและความรัดกุมทางทฤษฎี}

\begin{itemize}
    \item \textbf{แกนกลางปลอดการปนเปื้อน (Zero-Contamination Core):} ข้อมูลดิบทั้ง 40 รายการใน\\ \texttt{rhit\_v9\_0\_ai\_ssot.json} จำแนกออกเป็น \textbf{พารามิเตอร์พลวัตอิสระเชิงประสิทธิผล 18 ตัว} และ \textbf{พารามิเตอร์ที่ถูกตรึงด้วยสมมาตรและเฟส CP/Majorana 22 ตัว} ($\text{dof} = 4 > 0$) ค่าสังเกตการณ์พื้นฐานทั้ง 22 ค่าไม่ถูกเปลี่ยนแปลงใดๆ จากส่วนขยายทางดาราศาสตร์ฟิสิกส์ของรุ่น V9.0
    \item \textbf{การแยกระดับการพิสูจน์อย่างชัดเจน:} ในเสาหลักที่ 3 ค่าอันตรกิริยากำลังสี่ $\lambda_{\text{V9}}$ และรัศมีโซลิตอนสุญญากาศ $R_{\text{TF,vac}}$ เป็น \emph{การพิสูจน์เชิงเรขาคณิต 5 มิติโดยตรง (Exact 5D Geometric Derivations)} ขณะที่ความสัมพันธ์ระหว่างมวลฮาโลกับมวลดาวฤกษ์ ($\log_{10} M_{\text{floor}}, \alpha_{\text{SHMR}}, \log_{10} \rho_{\text{crit,tot}}$) เป็น \emph{สมมติฐานเชิงเรขาคณิตและปรากฏการณ์วิทยา (Geometric \& Phenomenological Ansatz)} ที่เทียบมาตรฐานกับข้อมูลจลนศาสตร์ LVDB
\end{itemize}

\section{เสาหลักที่ 1: ระบบฐานข้อมูลโฮโลแกรมและเพดานพลังงานตัดทอน UV}

เสาหลักที่ 1 กำหนดระดับพลังงานตัดทอนอัลตราไวโอเลต (UV) และความจุบิตข้อมูลบนขอบฟ้าจากสมการวิวัฒนาการกลุ่มการปรับมาตรฐาน (RGE) ใน 5 มิติ และกลไกเสถียรภาพ Goldberger-Wise (GW):

\begin{enumerate}
    \item \textbf{วิวัฒนาการ RGE สองระดับใน 5 มิติ:} คำนวณจากสเกลอิเล็กโทรวีค ($v_{\text{EW}} = 246.22\text{ GeV}$) ผ่านสเกลกลาง Pati-Salam ($M_I = 6.288608 \times 10^{12}\text{ GeV}$) สู่สเกลการรวมแรงครั้งใหญ่ ($M_{\text{GUT}} = 2.691598 \times 10^{16}\text{ GeV}$)
    \item \textbf{การนับจำนวนสถานะจำนวนเต็มในบัลก์ 5 มิติ:} การนับมัลติเพลตในบัลก์ ($n_{\mathbf{16}} = 0$, $n_{\mathbf{10}} = 2$, $n_{(15,1,1)} = 2$, $\sum T(R) = 22.0$) ให้ค่า $b_{B-L}^{5\text{D,eff}} = 8.0$ ซึ่งเมื่อรวมกับอัตราส่วนขีดเริ่มเปลี่ยน $r_T = 0.143202$ จะตรึงความลึกเชิงคอนฟอร์มัลอินฟราเรด:
    \begin{equation}
    u_{\text{IR}} = -\frac{8\pi^2}{b_{B-L}^{5\text{D,eff}} g_{\text{GUT}}^2} = -38.442470737632405
    \end{equation}
    \item \textbf{เพดานพลังงาน UV จำนวนบิตบนขอบฟ้า และอายุขัยโปรตอน:}
    \begin{align}
    E_{\text{UV}} &= v_{\text{EW}} e^{-u_{\text{IR}}} = 1.220889163672 \times 10^{19}\text{ GeV}, &
    l_{\text{cutoff}} &= \frac{\hbar}{M_{\text{UV}} c} = 1.616256302293 \times 10^{-35}\text{ m} \\
    N_{\text{bits}} &= \frac{4\pi R_H^2}{4 l_{\text{cutoff}}^2 \ln 2} = 3.252565292821 \times 10^{122}\text{ bits}, &
    \tau_p &\approx 1.51 \times 10^{48}\text{ ปี}
    \end{align}
\end{enumerate}

\section{เสาหลักที่ 2: แรงโน้มถ่วงอุบัติใหม่และพลวัตการพองตัวของเอกภพ}

ค่าคงที่ความโน้มถ่วงสากลของนิวตันอุบัติขึ้นโดยปราศจากการอ้างอิงวนซ้ำจากมวลตัดทอน $M_{\text{UV}}$ ร่วมกับการป้อนกลับของสุญญากาศอิเล็กโทรวีค:

\begin{equation}
G_N = \frac{\hbar c}{M_{\text{UV}}^2} \left[ 1 - \left(\frac{v_{\text{EW}}}{E_{\text{UV}}}\right)^4 \right] = 6.674310557973 \times 10^{-11}\text{ m}^3\text{kg}^{-1}\text{s}^{-2}
\end{equation}

ผลลัพธ์สอดคล้องกับค่ามาตรฐาน CODATA 2022 ($6.67430 \times 10^{-11}\text{ m}^3\text{kg}^{-1}\text{s}^{-2}$) ที่ $+0.0704\sigma$ ($\chi^2 = 0.0050$) สำหรับภาคการพองตัวแบบ Starobinsky $R^2$ ($N_e = 57.0212$ e-folds) ทฤษฎีทำนายอัตราส่วนคลื่นความโน้มถ่วงปฐมภูมิ $r = 12/N_e^2 = 0.00369070$ ($+0.0022\sigma$) และดัชนีสเปกตรัมสเกลาร์ $n_s = 1 - 2/N_e = 0.964925$ ($+0.0060\sigma$) พร้อมสมการวิวัฒนาการโฮโลแกรม:

\begin{equation}
G_N(z) = G_{N,0}\left[1 + \beta_{\text{holo}}\ln\left(\frac{H(z)}{H_0}\right)\right], \qquad \beta_{\text{holo}} = \frac{1}{\ln N_{\text{bits}}} \approx 0.00355
\end{equation}

\section{เสาหลักที่ 3: ภาคมืด แกนกลางโซลิตอน LVDB และเสียงสะท้อนของวัตถุความหนาแน่นสูงยิ่งยวด}

\subsection{เอกลักษณ์ทางพีชคณิตของพลังงานมืดโฮโลแกรม (OOA Reduced Form)}

ที่ค่าคงที่ฮับเบิล $H_0 = 67.5632\text{ km s}^{-1}\text{Mpc}^{-1}$ และสัดส่วนความพัวพันเชิงเรขาคณิต $\gamma_{\text{ent}} = 1 - 1/\pi \approx 0.681690$ การแบ่งส่วนพลังงานเชิงอุณหพลศาสตร์บนขอบฟ้า $S^2$ ให้เอกลักษณ์ทางพีชคณิตที่แม่นยำ:

\begin{equation}
\rho_{\text{DE}} = \frac{3\gamma_{\text{ent}} \hbar H_0^2}{8\pi l_{\text{cutoff}}^2 c} = \gamma_{\text{ent}} \left(\frac{3 H_0^2 c^2}{8\pi G_N}\right) \left[ 1 - \left(\frac{v_{\text{EW}}}{E_{\text{UV}}}\right)^4 \right] = 5.253185 \times 10^{-10}\text{ J m}^{-3}
\end{equation}

ค่านี้ตรงกับข้อมูลการสังเกตการณ์ของดาวเทียม Planck ($5.2532 \times 10^{-10}\text{ J m}^{-3}$) ที่ระดับ $-0.0002\sigma$ ($\chi^2 = 0.0000$)

\subsection{ส่วนขยาย V9.0 A: อันตรกิริยากำลังสี่ 5 มิติ และการเทียบมาตรฐานโซลิตอน LVDB}

ในรุ่น V8.5 สนามสเกลาร์อิสระ ($\lambda = 0$) ที่มีมวล $m_{\text{eff}} = 2.83 \times 10^{-20}\text{ eV}$ ให้แกนกลางควอนตัมขนาด $5.65\text{--}26.24\text{ pc}$ ซึ่งเล็กเกินไปสำหรับกาแล็กซีแคระทรงกลมคลาสสิก รุ่น RHIT V9.0 จึงเปิดใช้งานศักย์ผลักกำลังสี่ $V_{\text{sym}}(\phi) = -\frac{1}{2}\mu^2(\rho_{\text{tot}})\phi^2 + \frac{\lambda_{\text{V9}}}{4}\phi^4$ ผ่านเลขชี้กำลังเรขาคณิต 5 มิติ:

\begin{equation}
\lambda_{\text{V9}} = \exp\left[\left(4 + 2\delta_H - \frac{1}{2\pi}\right) u_{\text{IR}}\right] = 2.750175621636 \times 10^{-83}, \qquad
R_{\text{TF,vac}} = \pi \sqrt{\frac{\lambda_{\text{V9}} \hbar^3}{8\pi G_N m_{\text{eff}}^4 c}} = 320.3682\text{ pc}
\end{equation}

การบังคับใช้เกณฑ์แยกส่วนสุญญากาศ $\mu_0 \le (4\lambda_{\text{V9}}\rho_{\text{DE}}\cdot 10^{-15})^{1/4} \approx 1.290248 \times 10^{-27}\text{ eV}$ จำกัดการรบกวนต่อ $\rho_{\text{DE}}$ ให้ไม่เกิน $10^{-15}$ ขณะที่การพรางตัวในระบบสุริยะคงค่า $\langle\phi\rangle = 0$ ($\eta = 0.0$) ภายใต้การป้อนกลับของความหนาแน่น $-\frac{\phi^2}{2M_s^2}T^\mu_{\ \mu}$ รัศมีแกนกลางเชิงประสิทธิผล:

\begin{equation}
R_{\text{TF,eff}}(\rho_{\text{tot}}) = R_{\text{min}} + \frac{R_{\text{TF,vac}} - R_{\text{min}}}{1 + (\rho_{\text{tot}}/\rho_{\text{crit,tot}})^2}, \qquad (R_{\text{min}} = 12.0\text{ pc})
\end{equation}

จะเปลี่ยนผ่านเข้าสู่ชั้นห่อหุ้ม NFW ภายนอกอย่างราบรื่น:

\begin{center}
{\small
\setlength{\tabcolsep}{4.5pt}
\begin{tabular}{p{4.6cm} p{4.2cm} r r p{2.8cm}}
\toprule
\textbf{ชุดข้อมูล LVDB} & \textbf{แบบจำลองจลนศาสตร์} & \textbf{$\chi^2$ รวม} & \textbf{$\chi^2_{\text{red}}$} & \textbf{เทียบกับ NFW} \\
\midrule
บริวารทางช้างเผือก ($N = 46$) & มาตรฐาน $\Lambda\text{CDM}$ (NFW) & $1,355.62$ & $30.810$ & ฐานเปรียบเทียบ \\
\rowcolor{rowtint}
บริวารทางช้างเผือก ($N = 46$) & \textbf{RHIT V9.0 Hybrid ($T^\mu_{\ \mu}$ + TF)} & $\mathbf{804.44}$ & $\mathbf{18.708}$ & $\mathbf{-40.66\%}$ \\
\midrule
กลุ่มท้องถิ่นทั้งหมด ($N = 94$) & มาตรฐาน $\Lambda\text{CDM}$ (NFW) & $3,264.30$ & $35.482$ & ฐานเปรียบเทียบ \\
\rowcolor{rowtint}
กลุ่มท้องถิ่นทั้งหมด ($N = 94$) & \textbf{RHIT V9.0 Hybrid ($T^\mu_{\ \mu}$ + TF)} & $\mathbf{3,058.60}$ & $\mathbf{33.611}$ & $\mathbf{-6.30\%}$ \\
\bottomrule
\end{tabular}
}
\end{center}

\subsection{ส่วนขยาย V9.0 B: เสียงสะท้อนของวัตถุความหนาแน่นสูงยิ่งยวดและการสั่นพ้องของเมฆโบซอน}

ด้วยการส่งออกข้อมูลทางดาราศาสตร์ฟิสิกส์แบบทางเดียว (Zero-Contamination ต่อ SSOT) รุ่น RHIT V9.0 คำนวณมวลสูงสุดของโซลิตอนตามเกณฑ์ Kaup/Colpi ได้:

\begin{equation}
M_{\text{max,sol}} = \frac{\sqrt{\lambda_{\text{V9}}} M_{\text{UV}}^3}{m_{\text{eff}}^2} = 1.0683 \times 10^{16}\,M_\odot
\end{equation}

ค่าคงที่การจับคู่ของอะตอมความโน้มถ่วง $\alpha_g = G_N M m_{\text{eff}} / (\hbar c)$ และช่วงเวลาหน่วงเสียงสะท้อนของวัตถุไร้ขอบฟ้าเหตุการณ์ $\Delta t_{\text{echo}} = (2 r_s / c)|\ln \epsilon|$ ระหว่างขอบเขตตัดทอน UV ($\epsilon_{\text{UV}} = l_{\text{cutoff}}/r_s$) กับขอบเขตแกนกลางโซลิตอน V9.0 ($\epsilon_{\text{V9}} = \max[(r_s/R_{\text{TF,vac}})^2, \epsilon_{\text{UV}}]$):

\begin{center}
{\small
\setlength{\tabcolsep}{4.5pt}
\begin{tabular}{p{4.3cm} r r r r p{3.0cm}}
\toprule
\textbf{วัตถุเป้าหมายทางดาราศาสตร์} & \textbf{มวล ($M_\odot$)} & \textbf{$\alpha_g$} & \textbf{$\Delta t_{\text{V9}}$ (s)} & \textbf{$\Delta t_{\text{UV}}$ (s)} & \textbf{หอสังเกตการณ์} \\
\midrule
ซากหลุมดำ GW150914 & $62$ & $1.313 \times 10^{-8}$ & $0.0772$ & $0.1127$ & LIGO / Virgo / ET \\
\rowcolor{rowtint}
หลุมดำมวลปานกลาง (IMBH) & $1.0 \times 10^4$ & $2.118 \times 10^{-6}$ & $10.456$ & $19.172$ & LISA / ET \\
หลุมดำ Sgr A$^*$ (ทางช้างเผือก) & $4.297 \times 10^6$ & $9.100 \times 10^{-4}$ & $3,466.6$ & $8,751.6$ & EHT / ngEHT / LISA \\
\rowcolor{rowtint}
\textbf{หลุมดำ M87$^*$ (ใจกลาง Virgo A)} & $\mathbf{6.5 \times 10^9}$ & $\mathbf{1.377}$ & $\mathbf{3.369 \times 10^6}$ & $\mathbf{1.418 \times 10^7}$ & \textbf{EHT / ngEHT (สั่นพ้อง)} \\
\bottomrule
\end{tabular}
}
\end{center}

หลุมดำมวลยิ่งยวด M87$^*$ ($M = 6.5 \times 10^9\,M_\odot$) อยู่ในย่านการสั่นพ้องสูงสุดของปรากฏการณ์ Superradiance ($\alpha_g = 1.377 \in [0.2, 2.0]$) พอดี จึงเป็นเป้าหมายสำคัญในการทดสอบผ่านการถ่ายภาพโพลาไรเซชันและวงแหวนโฟตอนของกล้องโทรทรรศน์ ngEHT

\section{เสาหลักที่ 4: ภาคสสารและการแยกตัวประกอบเอกฐานเชิงซ้อน 5 มิติ}

มวลของเฟอร์มิออนมูลฐานทั้ง 12 ชนิด พารามิเตอร์การผสม CKM/PMNS ทั้ง 6 ค่า และค่าคงที่การละเมิดสมมาตร CP ของ Jarlskog ($J_{\text{CP}}$) อุบัติขึ้นจากสูตรการซ้อนทับใน 5 มิติเพียงสูตรเดียว โดยมีโปรไฟล์ฮิกส์ $\delta_H = \sqrt{4 + (m_H R_{\text{AdS}})^2} - 2 = 0.552141$:

\begin{equation}
Y_{ij}^{(f)} = \mathcal{I}_{5\text{D}}\!\left(c_{L,i}^{(f)},\, c_{R,j}^{(f)},\, u_{\text{IR}},\, \delta_H\right) \cdot y_{ij}^{(f)} e^{i\phi_{ij}^{(f)}}, \qquad f \in \{u, d, e, \nu\}
\end{equation}

มวลนิวทริโนเบาเกิดจากกลไก Type-I Seesaw $M_\nu = (M_D M_D^T)/M_I$ ให้ลำดับมวลปกติ $[3.251, 9.201, 50.216]\text{ meV}$

\section{ตารางตรวจสอบค่าสังเกตการณ์พื้นฐานครบทั้ง 22 ค่า (SSOT)}

{\small
\setlength{\tabcolsep}{4pt}
\begin{longtable}{p{2.7cm} l l p{4.1cm} r r}
\toprule
\textbf{หมวดหมู่} & \textbf{ตัวแปร} & \textbf{ค่าทำนาย} & \textbf{ค่าสังเกตการณ์ ($\pm 1\sigma$)} & \textbf{ส่วนเบี่ยงเบน} & \textbf{$\chi^2$} \\
\midrule
\endfirsthead
\toprule
\textbf{หมวดหมู่} & \textbf{ตัวแปร} & \textbf{ค่าทำนาย} & \textbf{ค่าสังเกตการณ์ ($\pm 1\sigma$)} & \textbf{ส่วนเบี่ยงเบน} & \textbf{$\chi^2$} \\
\midrule
\endhead
\bottomrule
\endfoot
จักรวาลวิทยาและแรงโน้มถ่วง & $G_N$ & $6.674311 \times 10^{-11}$ & $6.67430 \times 10^{-11} \pm 1.5 \times 10^{-15}$ & $+0.0704\sigma$ & $0.0050$ \\
 & $\rho_{\text{DE}}$ & $5.253185 \times 10^{-10}$ & $5.2532 \times 10^{-10} \pm 6.5 \times 10^{-12}$ & $-0.0002\sigma$ & $0.0000$ \\
 & $r_{\text{tensor}}$ & $0.003691$ & $0.00369 \pm 0.00032$ & $+0.0022\sigma$ & $0.0000$ \\
 & $n_s$ & $0.964925$ & $0.9649 \pm 0.0042$ & $+0.0060\sigma$ & $0.0000$ \\
\midrule
ควาร์ก & $m_t$ & $172.6861\text{ GeV}$ & $172.69 \pm 0.30\text{ GeV}$ & $-0.0129\sigma$ & $0.0002$ \\
 & $m_c$ & $1.2687\text{ GeV}$ & $1.270 \pm 0.020\text{ GeV}$ & $-0.0666\sigma$ & $0.0044$ \\
 & $m_u$ & $2.1342 \times 10^{-3}\text{ GeV}$ & $2.16 \times 10^{-3} \pm 4.9 \times 10^{-4}\text{ GeV}$ & $-0.0526\sigma$ & $0.0028$ \\
 & $m_b$ & $4.1839\text{ GeV}$ & $4.180 \pm 0.030\text{ GeV}$ & $+0.1316\sigma$ & $0.0173$ \\
 & $m_s$ & $9.0122 \times 10^{-2}\text{ GeV}$ & $9.34 \times 10^{-2} \pm 8.6 \times 10^{-3}\text{ GeV}$ & $-0.3812\sigma$ & $0.1453$ \\
 & $m_d$ & $4.6620 \times 10^{-3}\text{ GeV}$ & $4.67 \times 10^{-3} \pm 4.8 \times 10^{-4}\text{ GeV}$ & $-0.0167\sigma$ & $0.0003$ \\
\midrule
เลปตอนมีประจุ$^*$ & $m_\tau$ & $1.776856\text{ GeV}$ & $1.77686 \pm 0.0012\text{ GeV}$ & $-0.0032\sigma$ & $0.0000$ \\
 & $m_\mu$ & $0.105658\text{ GeV}$ & $0.105658 \pm 1.0 \times 10^{-6}\text{ GeV}$ & $+0.0232\sigma$ & $0.0005$ \\
 & $m_e$ & $5.1100 \times 10^{-4}\text{ GeV}$ & $5.110 \times 10^{-4} \pm 1.0 \times 10^{-6}\text{ GeV}$ & $-0.0017\sigma$ & $0.0000$ \\
\midrule
เมทริกซ์ CKM และการละเมิด CP & $V_{us}$ & $0.224346$ & $0.22430 \pm 0.00080$ & $+0.0578\sigma$ & $0.0033$ \\
 & $V_{cb}$ & $0.041643$ & $0.04130 \pm 0.00060$ & $+0.5713\sigma$ & $0.3264$ \\
 & $V_{ub}$ & $0.003988$ & $0.00382 \pm 0.00020$ & $+0.8382\sigma$ & $0.7025$ \\
 & $J_{\text{CP}}$ & $2.9515 \times 10^{-5}$ & $3.08 \times 10^{-5} \pm 1.5 \times 10^{-6}$ & $-0.8567\sigma$ & $0.7340$ \\
\midrule
นิวทริโนและเมทริกซ์ PMNS & $\Delta m_{31}^2$ & $2.5110 \times 10^{-3}\text{ eV}^2$ & $2.511 \times 10^{-3} \pm 2.7 \times 10^{-5}\text{ eV}^2$ & $+0.0018\sigma$ & $0.0000$ \\
 & $\Delta m_{21}^2$ & $7.4101 \times 10^{-5}\text{ eV}^2$ & $7.41 \times 10^{-5} \pm 2.1 \times 10^{-6}\text{ eV}^2$ & $+0.0005\sigma$ & $0.0000$ \\
 & $\sin^2\theta_{12}$ & $0.307000$ & $0.307 \pm 0.013$ & $-0.0000\sigma$ & $0.0000$ \\
 & $\sin^2\theta_{23}$ & $0.545983$ & $0.546 \pm 0.021$ & $-0.0008\sigma$ & $0.0000$ \\
 & $\sin^2\theta_{13}$ & $0.022030$ & $0.02203 \pm 0.00056$ & $+0.0003\sigma$ & $0.0000$ \\
\midrule
\multicolumn{5}{r}{\textbf{รวมทั้งหมด: 22 ค่าสังเกตการณ์ ($\text{dof} = 4$, สูงสุด $0.8567\sigma$) — Combined Total $\chi^2$}} & \textbf{1.9421} \\
\end{longtable}
}

\noindent{\footnotesize $^*$\textbf{หมายเหตุความคลาดเคลื่อนทางทฤษฎี:} ค่าความคลาดเคลื่อนขั้นต่ำ $1.0 \times 10^{-6}\text{ GeV}$ ($1\text{ keV}$) สำหรับเลปตอนมีประจุเบา ($m_e, m_\mu$) กำหนดขึ้นเพื่อรองรับการคำนวณระดับต้นไม้ (Tree-level 5D Yukawa matching) ก่อนรวมการแก้ไขเชิงรังสี QED/QCD ระดับ 2 ลูปจากสเกล $M_Z$ ลงมายัง $m_\ell$}

\section{แผนงานการทดสอบหักล้างเชิงการทดลอง (ค.ศ. 2026--2035)}

\begin{center}
{\small
\setlength{\tabcolsep}{4.5pt}
\begin{tabular}{p{2.2cm} p{3.5cm} p{6.1cm} p{3.2cm}}
\toprule
\textbf{เสาหลัก} & \textbf{ตัวชี้วัดการทดสอบ} & \textbf{เกณฑ์การหักล้างทฤษฎี} & \textbf{การทดลองเป้าหมาย} \\
\midrule
\textbf{เสาหลักที่ 1} & อายุขัยโปรตอน ($\tau_p$) & ทำนาย $\tau_p \approx 1.51 \times 10^{48}\text{ ปี}$ \textbf{ทฤษฎีจะถูกหักล้างหากพบการสลายตัวในช่วง $10^{34}\text{--}10^{36}\text{ ปี}$} & Hyper-Kamiokande, DUNE \\
\rowcolor{rowtint}
\textbf{เสาหลักที่ 2} & อัตราส่วนเทนเซอร์ ($r$) & ทำนาย $r = 0.00369070$ \textbf{ทฤษฎีจะถูกหักล้างหากวัดค่า $r \notin [0.00333, 0.00397]$} & Simons Observatory, LiteBIRD \\
\textbf{เสาหลักที่ 3-A} & แกนกลางโซลิตอนสสารมืด & ทำนาย $m_{\text{eff}} = 2.83 \times 10^{-20}\text{ eV}$, $R_{\text{TF,vac}} = 320.37\text{ pc}$ \textbf{ถูกหักล้างหากไม่พบแกนกลางโซลิตอนในกาแล็กซีแคระ} & JWST, Roman Space Telescope \\
\rowcolor{rowtint}
\textbf{เสาหลักที่ 3-B} & เมฆโบซอน M87$^*$ และเสียงสะท้อน GW & ทำนาย $\alpha_g = 1.377$ สำหรับ M87$^*$ \textbf{ถูกหักล้างหาก ngEHT ไม่พบเมฆโบซอนหรือไม่มีเสียงสะท้อนที่ $\ln B > 5$} & EHT / ngEHT, LIGO / ET / LISA \\
\textbf{เสาหลักที่ 4} & หลักความสมมูล ($\eta$) & กลไกพรางตัว Symmetron ในระบบสุริยะ ($\langle\phi\rangle = 0$) \textbf{ถูกหักล้างหากพบ $\eta > 10^{-15}$} & MICROSCOPE, Eot-Wash \\
\bottomrule
\end{tabular}
}
\end{center}

\end{document}